Atualmente, a implantação de conceitos de \gls{iot} nos mais diferentes dispositivos presentes 
no dia a dia de milhões de pessoas ao redor do mundo está numa crescente jamais antes vista. 
Como explicado por \citeonline{}, a adoção de dispositivos inteligentes trazem como benefício xxxx coisas. 

Em contrapartida, como explicado por \citeonline{fortune2025} e \citeonline{polimi2020}, 
o alto custo inicial exigido para realizar a transição de dispositivos analógicos para soluções inteligentes nativas,
 constitui o principal entrave para a adoção de tais tecnologias em larga escala. 
 A complexidade de infraestrutura necessária para a substituição física dos medidores limita a viabilidade econômica de projetos 
 em áreas privadas de pequeno e médio porte. 

Abordando o tema de automação residencial e comercial, 
o tema de smart metering surge como ponto focal na discussão. 
Como evidenciado pelo contraste em abordagens de \citeonline{} e \citeonline{}, 
o desenvolvimento de sistemas focados em smart metering 
enfrenta um impasse arquitetônico, 
onde o método de processamento divergia-se entre os projetos.

A implementação proposta por \citeonline{} é a denominada \gls{edge_computing}, 
onde o processamento é realizado localmente no dispositivo final, 
assim otimizando recursos energéticos e latência. 
Em oposição, a proposta de \citeonline{}, realizando processamento em cloud (\gls{cloud_computing}), 
surge como uma solução focada em facilidade de manutenção, 
utilização de softwares mais robustos e menor complexidade de implementação.

Neste contexto, a finalidade deste trabalho é realizar a comparação entre os métodos de processamento citados, 
definindo limites, vantagens e desvantagens sob determinado ambiente.
